拿 \(m=10^6\) 枚均匀硬币，每枚抛 \(n=100\) 次，挑出正面频率最低的那一枚，它的正面频率大约是 \(0.5-0.26\approx 0.24\)。一枚完全均匀的硬币，仅仅因为是从一百万个候选里被选出来的，就长出了偏币的外观。把每枚硬币看作一个候选假设、每次抛掷看作一个样本，同一句话就变成了：训练误差接近零的模型，在新数据上可能一塌糊涂。这不是工程事故，而是训练误差对真风险系统性乐观造成的数学后果。本单元要回答这份乐观有多大、由什么决定、又能怎样被控制住。

主线只有一条：把前面分开学过的三块材料接到同一个问题上。经验风险最小化（\hyperref[sec:09-Convex-Optimization/07-Applications/01]{``光滑经验风险与 \(L_2\) 正则化''}）负责把训练误差压下去；几何间隔（\hyperref[sec:09-Convex-Optimization/07-Applications/03]{``最大间隔与支持向量机''}）给出一个看上去更稳妥的选择准则；Hoeffding 集中界（\hyperref[sec:05-Probability/07-Transforms-and-Bounds/04]{``Chernoff 方法与 Hoeffding 集中界''}）提供量化``样本均值偏离期望多少''的尺子。四篇文章沿着一条推理链逐层推进。

\begin{center}
\begin{tikzpicture}[>=stealth]
\foreach \x/\lab in {0/{固定假设\\Hoeffding}, 3.1/{选择污染\\无偏性丢失}, 6.2/{联合界\\\(\log\abs{\mathcal{H}}\)}, 9.3/{VC 维\\生长函数}}
  {\draw (\x,0) rectangle (\x+2.5,1.1);
   \node[align=center,font=\footnotesize] at (\x+1.25,0.55) {\lab};}
\foreach \x in {2.5,5.6,8.7} {\draw[->] (\x,0.55) -- (\x+0.6,0.55);}
\draw[->] (10.55,0) -- (10.55,-0.7);
\draw (7.9,-1.8) rectangle (13.2,-0.7);
\node[align=center,font=\footnotesize] at (10.55,-1.25) {解释力与边界\\正则化、间隔、双下降};
\draw[dashed,gray] (1.25,0) -- (1.25,-1.25) -- (7.9,-1.25);
\node[gray,font=\footnotesize,anchor=north] at (4.2,-1.3) {测试集只能用一次};
\end{tikzpicture}
\end{center}

\emph{四篇的推理链：从一个干净的固定假设出发，被数据依赖的选择打断，用联合界修复有限类，再用 VC 维接管无限类，最后盘点这条链能承诺什么。}

\begin{enumerate}
\tightlist
\item
  训练误差为何总是乐观：定义真风险与经验风险，说明同一批数据被``选假设''与``评假设''用了两次会让训练误差向下偏，以及单一固定假设的 Hoeffding 界为什么对选出来的 \(\widehat h\) 失效。
\item
  从联合界到一致收敛：假设类有限时，联合界以 \(\log\abs{\mathcal{H}}\) 的代价换来对所有假设同时成立的界，ERM 因此获得超出风险控制；同时看到任何带实值参数的假设类都会让这条路直接断掉。
\item
  VC 维把无限假设类变有限：用``在 \(n\) 个点上能产生多少种标记''替代类的基数，经生长函数、打散与 Sauer 引理得到不依赖 \(\abs{\mathcal{H}}\) 的泛化界，并给出``需要多少样本''的量级答案。
\item
  泛化界的解释力与边界：承认这些界在数值上很松，同时看清它们如何解释正则化、结构风险最小化与间隔，并直面过参数化与双下降这些经典框架尚未覆盖的现代现象。
\end{enumerate}

\subsection{怎么读这一章}

随手翻阅时，可以只看第一篇的硬币例子与``选择和评估为何不能共用数据''，第二篇的``同时控制为什么要为候选数量付费''，第三篇的``容量是可实现的行为而非参数个数''，第四篇的``界是结构判断而非性能预言''。想弄清证明机制，再回头读 Hoeffding、联合界、ERM 的三步不等式、Sauer 引理与对称化骨架。真正要把结论用在自己的项目上时，还得逐项核对独立同分布、损失有界、假设类是否预先确定，以及数据是否已经参与过模型选择——本单元所有的界，条件一旦松动就全部作废。
